\newvideofile{einfuehrungElektrischeMaschinen}{Einführung}

\s{
    \section{Klassifizierung elektrischer Maschinen}
    \glqq Elektrische Maschine\grqq\, ist eine andere Bezeichnung für einen elektromagnetische Energiewandler und beschreibt zum einen den Transformator als ruhende elektrische Maschine und zum anderen einen Elektromotore oder Generator.
    Elektromotoren treiben viele alltägliche Geräte wie Haushaltsgeräte und Elektrofahrzeuge an, während Generatoren und Transformatoren die Erzeugung und Verteilung von elektrischer Energie sicherstellen.

    Dieses Kapitel bietet eine Einführung in die grundlegenden Konzepte und Funktionsweisen elektrischer Maschinen. Zunächst werden die verschiedenen Typen elektrischer Maschinen vorgestellt, einschließlich ihrer Prinzipien der Energieumwandlung und ihrer Hauptkomponenten. Es folgt eine detaillierte Betrachtung der Funktionsweise von Elektromotoren und Generatoren, der Relevanz von Transformatoren für die Energieübertragung. Außerdem soll der Einfluss neuer Technologien auf die Effizienz und Leistung elektrischer Maschinen betrachtet werden.

    \begin{Lernziele}{Elektrische Maschinen}
        Die Studierenden
        \begin{itemize}
            \item kennen die verschiedenen elektrischen Maschinen.
            \item können grundlegende Aufgabenstellung im Themenbereich der elektrischen Maschinen lösen.
        \end{itemize}
    \end{Lernziele}
}
\b{
    \begin{frame}{Lernziele}
        \begin{Lernziele}{Elektrische Maschinen}
            Die Studierenden
            \begin{itemize}
                \item kennen die verschiedenen Arten elektrischer Maschinen.
                \item können grundlegende Aufgabenstellung im Themenbereich der elektrischen Maschinen lösen.
            \end{itemize}
        \end{Lernziele}
    \end{frame}
}
\v{
    %Videoframe
    \begin{frame}{Elektrische Maschinen}
        \speech{einfuehrungElektrischeMaschinen1}{1}{Der Begriff, Elektrische Maschinen, ist eine andere Bezeichnung für elektromagnetische Energiewandler und beschreibt zum einen den Transformator als ruhende elektrische Maschine und zum anderen Elektromotoren und Generatoren.}
        \silence{2}
    \end{frame}
    \begin{frame}{Elektrische Maschinen}
        \f{}{width=0.3\textwidth}{Transformator2b.png}{}
        \speech{einfuehrungElektrischeMaschinen2}{1}{Elektrische Maschinen sind Energieumwandler. Es wird vornehmlich zwischen Transformatoren und Motoren bzw. Generatoren unterschieden. Transformatoren wandeln elektrische Energie in eine andere Form von elektrischer Energie um, indem sie ein Spannungsniveau auf ein anderes übertragen.} %### neuer Seed
        \silence{2}
    \end{frame}
    \begin{frame}{Elektrische Maschinen}
        \f{}{width=0.5\textwidth}{Aufbau_Synchronmaschine1}{}
        \speech{einfuehrungElektrischeMaschinen3}{1}{Motoren und Generatoren hingegen wandeln elektrische Energie in mechanische Energie um und umgekehrt. Motoren konvertieren elektrische Energie in mechanische Energie, während Generatoren mechanische Energie in elektrische Energie umwandeln. Da der Umwandlungsprozess umkehrbar ist, können Motoren auch als Generatoren und umgekehrt eingesetzt werden. Eine weitere Klassifizierung bei Motoren bzw. Generatoren wird je nach betriebenen Strom in Gleichstrommaschinen und Wechselstromaschinen eingeteilt. Die Wechselstrommaschinen werden zudem in Synchron- und Asychronmaschinen gegliedert. Die Funktionsweisen aller elektrischen Maschinen basieren, auf den in Modul 6 besprochenen, elektromagnetischen Prinzipien der Induktion und der Lorentz-kraft. Während der Transformator rein auf dem Prinzip der Induktion beruht, nutzt der Motor beziehungsweise der Generator je nach Wandlungsprozess unterschiedliche elektromagnetische Prinzipien. In der Funktion als Generator basiert er auf Induktion, um elektrische Energie durch sich ändernde Magnetfelder zu erzeugen. Das Prinzip der Lorentz-kraft macht er sich in seiner Funktion als als Motor zunutze.}
        \silence{2}
    \end{frame}
    \begin{frame}{Lernziele}
        \begin{Lernziele}{Elektrische Maschinen}
            Du kannst
            \begin{itemize}
                \item die verschiedenen Arten elektrischer Maschinen benennen und voneinander unterscheiden.
                \item grundlegende Aufgabenstellung im Themenbereich der elektrischen Maschinen lösen.
            \end{itemize}
        \end{Lernziele}
        \speech{einfuehrungElektrischeMaschinen4}{1}{Dieses Modul bietet mit den nachfolgenden Videos eine Einführung in die grundlegenden Konzepte und Funktionsweisen elektrischer Maschinen. Dazu scheun wir die zunächst die verschiedenen Typen elektrischer Maschinen an, einschließlich ihrer Prinzipien der Energieumwandlung und ihrer Hauptkomponenten. Danach werfen wir einen detaillierten Blick auf die Funktionsweise von Elektromotoren und Generatoren und auf die Relevanz von Transformatoren für die Energieübertragung. Wir sehen uns in diesem Modul außerdem den Einfluss neuer Technologien auf die Effizienz und Leistung elektrischer Maschinen an.}
        \silence{2}
    \end{frame}
}
\s{
    Elektrische Maschinen sind Energieumwandler. Es wird vornehmlich zwischen Transformatoren und Motoren bzw. Generatoren unterschieden. Transformatoren wandeln elektrische Energie in eine andere Form von elektrischer Energie um, indem sie ein Spannungsniveau auf ein anderes übertragen.
    Motoren und Generatoren hingegen wandeln elektrische Energie in mechanische Energie um und umgekehrt. Motoren konvertieren elektrische Energie in mechanische Energie, während Generatoren mechanische Energie in elektrische Energie umwandeln. Da der Umwandlungsprozess umkehrbar ist, können Motoren auch als Generatoren und umgekehrt eingesetzt werden. Eine weitere Klassifizierung bei Motoren bzw. Generatoren findet je nach betriebenen Strom in Gleichstrommaschinen und Wechselstromaschinen statt. Die Wechselstrommaschinen werden des Weiteren in Synchron- und Asynchronmaschinen gegliedert.


    Die Funktionsweisen aller elektrischen Maschinen basieren auf den in Modul 6 besprochenen elektromagnetischen Prinzipien der Induktion und der Lorentzkraft.
    Während der Transformator rein auf dem Prinzip der Induktion beruht, nutzt der Motor bzw. Generator je nach Wandlungsprozess unterschiedliche elektromagnetische Prinzipien. In der Funktion als Generator basiert er auf Induktion, um elektrische Energie durch sich ändernde Magnetfelder zu erzeugen. Das Prinzip der Lorentzkraft macht er sich als Motor zunutzen.
}
